% Values: PRECISE, FAITHFUL
% Evidence planned:
% - PRECISE: B0 = ME0 and E0 given module by module, every interface as a
%   triple (RA, CE, VE) of Chapter 7, the dynamic operator table OT_DYN, every
%   identifier of the initial static basis (Appendix C) with its value, every
%   basic value with its arity; APPLY defined for every basic value by displayed definitions:
%   the floats built from their significands and exponents and rounded by fl;
%   Int arithmetic through the parameters IntRange and norm; IEEE binary64
%   arithmetic case by case; powers and elementary functions with their
%   special values; conversions with saturation; equality (Eq) and order
%   (cmp) by structural recursion; min, max and clamp (A48, A49);
%   Booleans, append, Debug, arrays and
%   byte sequences (A92); the
%   values written in Elm given by their elm/core 1.0.5 text and evaluated
%   by the rules of Chapter 6; effect values by reference to Chapter 8.
% - FAITHFUL: D3 (every operation total: 3 // 0 is 0), D4 (the basic
%   mathematics defined here, not in a library), D8 (IntRange, norm,
%   characters, accuracy and Debug.toString left to Appendix E); rulings
%   A40 to A65 cited beside the definitions they shape; I2: the same
%   identifiers as Appendix C, module by module.

\chapter{The Initial Dynamic Basis}\label{ch:appendix-dynamic-basis}

We indicate components of the initial basis by the subscript 0, as in
Appendix~\ref{ch:appendix-static-basis}; in this appendix $\ME_0$, $E_0$ and
$\VE_0$ denote dynamic objects. The initial dynamic basis is the module
environment
\[
  B_0=\ME_0 ,
\]
whose domain is the domain of the static $\ME_0$
(Section~\ref{sec:appendix-static-basis-modenv}). For each $\modid$ in that
domain, $\ME_0(\modid)=(\dmRA,\dmCE,\VE)$ is a dynamic interface
(Section~\ref{sec:dynamic-modules-compound}), given in
Section~\ref{sec:appendix-dynamic-basis-env}. Every module is evaluated
starting from the environment
\[
  E_0=(\emptymap,\VE_0),\qquad
  \VE_0=\{\kw{[]}\mapsto\kw{[]},\ \kw{::}\mapsto\kw{::}\},
\]
to which its imports add (Chapter~\ref{ch:dynamic-modules}), as statically it
starts from the environment $E_0$ of
Section~\ref{sec:appendix-static-basis-modenv}. $E_0$ is an environment of
Chapter~\ref{ch:dynamic-core}, with an empty qualified environment; it binds
the two list constructor names to the constructor values they name
(Section~\ref{sec:dynamic-core-simple}). % decision 2 (audit X4)
Every identifier of the initial static basis is bound in the initial dynamic
basis, and the reverse; the hidden identifiers of
Section~\ref{sec:appendix-static-basis-modenv} are bound only in the
environments of the texts of Section~\ref{sec:appendix-dynamic-basis-elm}. % I2

The values of the basis are constructor values, special values, closures
given by the text of the library \kw{elm/core} 1.0.5, and basic values, whose
meaning is the function $\APPLY$ of Section~\ref{sec:appendix-dynamic-basis-apply}.
The basic mathematical operations are defined there in full, on the special
values defined in Section~\ref{sec:appendix-dynamic-basis-svals}; every one
of them has a result for every argument, a defined value being given even
where no mathematically correct one exists. % D3, D4
Three things are left to the implementation and recorded in
Appendix~\ref{ch:appendix-implementation}: the range of Int with its
normalisation, the accuracy of the elementary functions on floats, and the
string that \kw{Debug.toString} gives. % D8

% ===========================================================================
\section{Special Values}\label{sec:appendix-dynamic-basis-svals}

The class of special values (Section~\ref{sec:dynamic-core-simple}) is the
disjoint union
\[
  \SVal=\IntRange\cup\effFlt\cup\dbChr\cup\effStr
\]
of the \emph{Int values}, the \emph{floats}, the \emph{characters} and the
\emph{strings}, defined below. The union is disjoint: the Int value $3$ and
the float $3$ are different special values. We write $\mathbb{Z}$ for the set
of integers and $\mathbb{R}$ for the set of real numbers.

\subsection{Int Values}\label{sec:appendix-dynamic-basis-ints}

The set $\IntRange\subseteq\mathbb{Z}$ of Int values and the
\emph{normalisation} $\norm:\mathbb{Z}\to\IntRange$ are parameters of the
definition. Every implementation chooses them so that:
% A40, A45
\begin{itemize}
\item $\IntRange$ has a least member $\min\IntRange$ and a greatest member
  $\max\IntRange$, and contains every integer between them;
\item $\min\IntRange\leq -2^{31}$ and $2^{31}-1\leq\max\IntRange$;
\item if $n\in\IntRange$ and $n\neq\min\IntRange$, then $-n\in\IntRange$;
\item $\norm(n)=n$ for every $n\in\IntRange$.
\end{itemize}
% DECIDED: IntRange is an interval with a least and a greatest member; the
% saturating conversions of A44 need both bounds.
The choices of the implementations are given in
Section~\ref{sec:appendix-implementation-int}. Int values are integers, so Int
has one zero. % A45

\subsection{Floats}\label{sec:appendix-dynamic-basis-floats}

The \emph{finite non-zero floats} are the real numbers $m\cdot 2^{q}$ and
$-m\cdot 2^{q}$, where $m$ and $q$ are integers with either
\[
  2^{52}\leq m<2^{53}\ \text{and}\ -1074\leq q\leq 971,
  \quad\text{or}\quad
  1\leq m<2^{52}\ \text{and}\ q=-1074 .
\]
Each has exactly one such representation, and $m$ is its \emph{significand}.
They are the finite non-zero numbers of the binary64 format of the standard
IEEE 754-2019. The set of floats is
\[
  \effFlt=\{\text{finite non-zero floats}\}\cup\{+0,-0,+\infty,-\infty,\dbNaN\},
\]
and $\effFltFin=\effFlt\setminus\{+\infty,-\infty,\dbNaN\}$ is the set of
\emph{finite} floats. The two zeros are different floats, and so are the two
infinities. There is one NaN.
% DECIDED: the many NaNs of the binary64 format are one float here: no basic
% value of the initial basis distinguishes them.
In this appendix $x$, $y$ and $z$ range over $\effFlt$.

In a formula, a float other than NaN stands for the number it denotes on the
extended real line, both zeros denoting $0$; thus $x<y$ and $x=y$ compare
numbers, and $+0=-0$ holds as a comparison of numbers. That $x$ \emph{is} a
particular float, such as $-0$, is said in words. For $x\neq\dbNaN$, the
\emph{sign} $\dbsgn(x)$ is $+1$ when $x$ is $+0$, $+\infty$ or positive, and
$-1$ when $x$ is $-0$, $-\infty$ or negative. For $s\in\{+1,-1\}$, $s\cdot 0$
and $s\cdot\infty$ are the zero and the infinity of sign $s$, and for a float
$z\neq\dbNaN$, $s\cdot z$ is $z$ when $s=+1$ and $z$ with its sign reversed
when $s=-1$.

Real numbers are rounded to floats by $\dbfl:\mathbb{R}\to\effFlt$, rounding
to nearest with ties to even as IEEE 754-2019 defines it:
\begin{itemize}
\item $\dbfl(r)=+\infty$ if $r\geq 2^{1024}-2^{970}$, and
  $\dbfl(r)=-\infty$ if $r\leq -(2^{1024}-2^{970})$;
\item otherwise, if $r\neq 0$, $\dbfl(r)$ is the finite float nearest to $r$
  and, of two equally near, the one whose significand is even (a zero counting
  as even); when that float is a zero, it is the zero of the sign of $r$;
\item $\dbfl(0)=+0$.
\end{itemize}
Every Int value $n$ with $|n|\leq 2^{53}$ is a float exactly:
$\dbfl(n)=n$.

\subsection{Characters and Strings}\label{sec:appendix-dynamic-basis-chars}

The set $\dbChr$ of characters is a parameter of the definition
(Section~\ref{sec:appendix-implementation-chars}), with an injective function
\[
  \dbcode:\dbChr\to\IntRange,
\]
the \emph{code} of a character, which the basic value \kw{Char.toCode} gives
(Figure~\ref{fig:dynamic-basis-conv}). For every Unicode scalar value $h$
there is a character with the code $h$, the character with code point $h$:
it is the character that the escape \kw{\textbackslash u\kwlb}$h$\kw{\kwrb}
denotes (Figure~\ref{fig:core-escapes}), and the one that a character
constant holding that character itself denotes. A further parameter is a
function
\[
  \mathit{other}:\mathbb{Z}\to\dbChr,
\]
the character that \kw{Char.fromCode} gives for an integer that is not a
Unicode scalar value. Nothing is required of $\mathit{other}$: its result may
be a character that is no Unicode scalar value, such as a lone surrogate, but
it is a character, and so has a code in $\IntRange$ distinct from that of
every other character. Nothing else is required of characters: the initial
basis has no operation on characters other than equality, order,
\kw{Char.toCode} and \kw{Char.fromCode}. % A52
A string is a finite sequence $\dbch_1\cdots\dbch_k$ of characters
($k\geq 0$), its \emph{length} being $k$; $\effStr$ is the set of strings, and
$s$ ranges over $\effStr$. A character constant denotes a character and a
string constant a string (Chapter~\ref{ch:syntax-core}). % A53

% ===========================================================================
\section{The Value Environments}\label{sec:appendix-dynamic-basis-env}

Every interface of $\ME_0$ has $\dmRA=\emptyset$, since the basis has no
record alias. The constructor environments are
\[
\begin{array}{@{}l@{\ =\ }l@{}}
  \dmCE\of\ME_0(\kw{Basics}) &
    \{\kw{Bool}\mapsto\{\kw{True}\mapsto\kw{True},\ \kw{False}\mapsto\kw{False}\},\\
  \multicolumn{1}{l}{} & \phantom{\{}\kw{Order}\mapsto\{\kw{LT}\mapsto\kw{LT},\
    \kw{EQ}\mapsto\kw{EQ},\ \kw{GT}\mapsto\kw{GT}\},\\
  \multicolumn{1}{l}{} & \phantom{\{}\kw{Never}\mapsto\emptymap\},\\
  \dmCE\of\ME_0(\kw{Maybe}) &
    \{\kw{Maybe}\mapsto\{\kw{Just}\mapsto\kw{Just},\ \kw{Nothing}\mapsto\kw{Nothing}\}\},\\
  \dmCE\of\ME_0(\kw{Result}) &
    \{\kw{Result}\mapsto\{\kw{Ok}\mapsto\kw{Ok},\ \kw{Err}\mapsto\kw{Err}\}\},
\end{array}
\]
and $\dmCE\of\ME_0(\modid)=\emptymap$ for every other $\modid$. Thus $\dmCE$
maps every custom type that a module of the basis exposes to the constructors
it exposes, the closed type \kw{Never} to the empty map; the abstract types
are in the domain of no $\dmCE$. % decision 3 (audit X2, Y2)
The value environments of the interfaces of $\ME_0$ are shown in
Figures~\ref{fig:dynamic-basis-basics} to~\ref{fig:dynamic-basis-others}.
In the figures, a constructor is bound to the constructor value it names,
whose arity is shown. An identifier shown with an arity is bound to the basic
value of the same name, which has that arity; we write the basic values of
\kw{Basics} unqualified (\kw{add}) and those of other modules qualified
(\kw{Debug.log}), and the basic values of
Figures~\ref{fig:dynamic-basis-basics} to~\ref{fig:dynamic-basis-others} are
all distinct. An identifier marked ``Elm'' is bound to the value of its
elm/core text (Section~\ref{sec:appendix-dynamic-basis-elm}); for \kw{abs},
to that of its default copy, the copy at $\tFloat$ being bound to a new
variable (Section~\ref{sec:appendix-dynamic-basis-elm}). % A49
The identifiers
\kw{e} and \kw{pi} are bound to floats.

\begin{figure}[htbp]
\centering
\small
\[
\begin{array}{|r@{\ }c@{\ }l@{\qquad}c|}
\hline
\var\ \text{or}\ \con & \mapsto & v & \arity\\
\hline
\kw{True}\citealso{b:Basics.Bool},\ \kw{False}\citealso{b:Basics.Bool} & \mapsto & \kw{True},\ \kw{False} & 0\\
\kw{LT}\citealso{b:Basics.Order},\ \kw{EQ}\citealso{b:Basics.Order},\ \kw{GT}\citealso{b:Basics.Order} & \mapsto & \kw{LT},\ \kw{EQ},\ \kw{GT} & 0\\
\hline
\kw{toFloat}\citealso{b:Basics.toFloat},\ \kw{round}\citealso{b:Basics.round},\ \kw{floor}\citealso{b:Basics.floor},\ \kw{ceiling}\citealso{b:Basics.ceiling},\ \kw{truncate}\citealso{b:Basics.truncate} & \mapsto & \text{basic value} & 1\\
\kw{compare}\citealso{b:Basics.compare} & \mapsto & \text{basic value} & 2\\
\kw{min}\citealso{b:Basics.min},\ \kw{max}\citealso{b:Basics.max} & \mapsto & \text{basic value} & 2\\
\kw{clamp}\citealso{b:Basics.clamp} & \mapsto & \text{basic value} & 3\\
\kw{not}\citealso{b:Basics.not} & \mapsto & \text{basic value} & 1\\
\kw{xor}\citealso{b:Basics.xor} & \mapsto & \text{basic value} & 2\\
\kw{modBy}\citealso{b:Basics.modBy},\ \kw{remainderBy}\citealso{b:Basics.remainderBy} & \mapsto & \text{basic value} & 2\\
\kw{negate}\citealso{b:Basics.negate} & \mapsto & \text{basic value} & 1\\
\kw{sqrt}\citealso{b:Basics.sqrt} & \mapsto & \text{basic value} & 1\\
\kw{cos}\citealso{b:Basics.cos},\ \kw{sin}\citealso{b:Basics.sin},\ \kw{tan}\citealso{b:Basics.tan},\ \kw{acos}\citealso{b:Basics.acos},\ \kw{asin}\citealso{b:Basics.asin},\ \kw{atan}\citealso{b:Basics.atan} & \mapsto & \text{basic value} & 1\\
\kw{atan2}\citealso{b:Basics.atan2} & \mapsto & \text{basic value} & 2\\
\kw{isNaN}\citealso{b:Basics.isNaN},\ \kw{isInfinite}\citealso{b:Basics.isInfinite} & \mapsto & \text{basic value} & 1\\
\hline
\kw{e}\citealso{b:Basics.e} & \mapsto & \dbfl(\dbeuler) & \\
\kw{pi}\citealso{b:Basics.pi} & \mapsto & \dbfl(\dbpinum) & \\
\hline
\kw{abs}\citealso{b:Basics.abs} & \mapsto & \text{Elm} & \\
\kw{logBase}\citealso{b:Basics.logBase},\ \kw{degrees}\citealso{b:Basics.degrees},\ \kw{radians}\citealso{b:Basics.radians},\ \kw{turns}\citealso{b:Basics.turns} & \mapsto & \text{Elm} & \\
\kw{fromPolar}\citealso{b:Basics.fromPolar},\ \kw{toPolar}\citealso{b:Basics.toPolar} & \mapsto & \text{Elm} & \\
\kw{identity}\citealso{b:Basics.identity},\ \kw{always}\citealso{b:Basics.always},\ \kw{never}\citealso{b:Basics.never} & \mapsto & \text{Elm} & \\
\hline
\end{array}
\]
\caption{Dynamic $\VE$ of $\ME_0(\kw{Basics})$}\label{fig:dynamic-basis-basics}
\end{figure}

Each operator of the initial basis stands for a hidden variable of
\kw{Basics} or, for \kw{::}, the list constructor
(Figure~\ref{fig:static-basis-operators}). The \emph{dynamic operator table}
\[
  \OTdyn : \cls{Op} \finmap \Val
\]
maps each operator of Figure~\ref{fig:core-fixities} to the value of the
identifier it stands for, and is given in
Figure~\ref{fig:dynamic-basis-operators}. The value of an operator is
$\OTdyn(\opv)$ in every environment: it is looked up in no environment and
reached through no interface. % A16, A56; decision 1 (audit X1)

\begin{figure}[htbp]
\centering
\small
\[
\begin{array}{|c@{\quad}l@{\ }c@{\ }l@{\qquad}c|}
\hline
\opv & \text{stands for} & \mapsto & \OTdyn(\opv) & \arity\\
\hline
\kw{<<},\ \kw{>>} & \kw{composeL}\citealso{b:Basics.composeL},\ \kw{composeR}\citealso{b:Basics.composeR} & \mapsto & \text{Elm} & \\
\kw{\textasciicircum} & \kw{pow}\citealso{b:Basics.pow} & \mapsto & \text{basic value} & 2\\
\kw{*},\ \kw{/},\ \kw{//} & \kw{mul}\citealso{b:Basics.mul},\ \kw{fdiv}\citealso{b:Basics.fdiv},\ \kw{idiv}\citealso{b:Basics.idiv} & \mapsto & \text{basic value} & 2\\
\kw{+},\ \kw{-} & \kw{add}\citealso{b:Basics.add},\ \kw{sub}\citealso{b:Basics.sub} & \mapsto & \text{basic value} & 2\\
\kw{++} & \kw{append}\citealso{b:Basics.append} & \mapsto & \text{basic value} & 2\\
\kw{::} & \kw{::}\citealso{b:List} & \mapsto & \kw{::} & 2\\
\kw{==},\ \kw{/=} & \kw{eq}\citealso{b:Basics.eq},\ \kw{neq}\citealso{b:Basics.neq} & \mapsto & \text{basic value} & 2\\
\kw{<},\ \kw{>},\ \kw{<=},\ \kw{>=} & \kw{lt}\citealso{b:Basics.lt},\ \kw{gt}\citealso{b:Basics.gt},\ \kw{le}\citealso{b:Basics.le},\ \kw{ge}\citealso{b:Basics.ge} & \mapsto & \text{basic value} & 2\\
\kw{\&\&},\ \kw{||} & \kw{and}\citealso{b:Basics.and},\ \kw{or}\citealso{b:Basics.or} & \mapsto & \text{basic value} & 2\\
\kw{|>},\ \kw{<|} & \kw{apR}\citealso{b:Basics.apR},\ \kw{apL}\citealso{b:Basics.apL} & \mapsto & \text{Elm} & \\
\hline
\end{array}
\]
\caption{The Dynamic Operator Table, $\OTdyn$}\label{fig:dynamic-basis-operators}
\end{figure}

\begin{figure}[htbp]
\centering
\small
\[
\begin{array}{|r@{\ }c@{\ }l@{\qquad}c|}
\hline
\var\ \text{or}\ \con & \mapsto & v & \arity\\
\hline
\multicolumn{4}{|l|}{\kw{Maybe}}\\
\hline
\kw{Just}\citealso{b:Maybe.Maybe} & \mapsto & \kw{Just} & 1\\
\kw{Nothing}\citealso{b:Maybe.Maybe} & \mapsto & \kw{Nothing} & 0\\
\hline
\multicolumn{4}{|l|}{\kw{Result}}\\
\hline
\kw{Ok}\citealso{b:Result.Result},\ \kw{Err}\citealso{b:Result.Result} & \mapsto & \kw{Ok},\ \kw{Err} & 1\\
\hline
\multicolumn{4}{|l|}{\kw{Char}}\\
\hline
\kw{toCode}\citealso{b:Char.toCode},\ \kw{fromCode}\citealso{b:Char.fromCode} & \mapsto & \text{basic value} & 1\\
\hline
\multicolumn{4}{|l|}{\kw{Debug}}\\
\hline
\kw{log}\citealso{b:Debug.log} & \mapsto & \text{basic value} & 2\\
\kw{todo}\citealso{b:Debug.todo},\ \kw{toString}\citealso{b:Debug.toString} & \mapsto & \text{basic value} & 1\\
\hline
\multicolumn{4}{|l|}{\kw{Platform}}\\
\hline
\kw{worker}\citealso{b:Platform.worker} & \mapsto & \text{basic value} & 1\\
\hline
\multicolumn{4}{|l|}{\kw{Platform.Cmd}\ \text{and}\ \kw{Platform.Sub}}\\
\hline
\kw{none}\citealso{b:Platform.Cmd.none}\citealso{b:Platform.Sub.none} & \mapsto & \text{Elm} & \\
\kw{batch}\citealso{b:Platform.Cmd.batch}\citealso{b:Platform.Sub.batch} & \mapsto & \text{basic value} & 1\\
\kw{map}\citealso{b:Platform.Cmd.map}\citealso{b:Platform.Sub.map} & \mapsto & \text{basic value} & 2\\
\hline
\multicolumn{4}{|l|}{\kw{Task}}\\
\hline
\kw{succeed}\citealso{b:Task.succeed},\ \kw{fail}\citealso{b:Task.fail} & \mapsto & \text{basic value} & 1\\
\kw{andThen}\citealso{b:Task.andThen},\ \kw{onError}\citealso{b:Task.onError} & \mapsto & \text{basic value} & 2\\
\kw{map}\citealso{b:Task.map},\ \kw{map2}\citealso{b:Task.map2},\ \kw{map3}\citealso{b:Task.map3},\ \kw{map4}\citealso{b:Task.map4},\ \kw{map5}\citealso{b:Task.map5} & \mapsto & \text{Elm} & \\
\kw{sequence}\citealso{b:Task.sequence},\ \kw{mapError}\citealso{b:Task.mapError},\ \kw{perform}\citealso{b:Task.perform},\ \kw{attempt}\citealso{b:Task.attempt} & \mapsto & \text{Elm} & \\
\hline
\multicolumn{4}{|l|}{\kw{Process}}\\
\hline
\kw{spawn}\citealso{b:Process.spawn},\ \kw{sleep}\citealso{b:Process.sleep},\ \kw{kill}\citealso{b:Process.kill} & \mapsto & \text{basic value} & 1\\
\hline
\end{array}
\]
\caption{Dynamic $\VE$ of the Other Modules of $\ME_0$}\label{fig:dynamic-basis-others}
\end{figure}

The interfaces of the modules \kw{List}, \kw{Tuple}, \kw{String},
\kw{Array}, \kw{Bytes}, \kw{Json.\allowbreak Encode}, \kw{Json.\allowbreak Decode},
\kw{VirtualDom} and \kw{Html} are $(\emptyset,\emptymap,\emptymap)$: these
modules are in $\ME_0$ for their types, which have no dynamic counterpart
(the values of \kw{Array} and \kw{Bytes} being given in
Sections~\ref{sec:appendix-dynamic-basis-arrays}
and~\ref{sec:appendix-dynamic-basis-bytes}). The interfaces of \kw{Char}, \kw{Platform},
\kw{Platform.Cmd}, \kw{Platform.Sub}, \kw{Task}, \kw{Process} and \kw{Debug}
are $(\emptyset,\emptymap,\VE)$, with $\VE$ as in
Figure~\ref{fig:dynamic-basis-others}. These interfaces have the property of
Section~\ref{sec:dynamic-modules-compound}. The basic values
\kw{Platform.\allowbreak Cmd.\allowbreak batch} and
\kw{Platform.\allowbreak Cmd.\allowbreak map}, and their counterparts in \kw{Platform.Sub}, are written
\kw{Cmd.batch}, \kw{Cmd.map}, \kw{Sub.batch} and \kw{Sub.map} in
Chapter~\ref{ch:effects}, after the qualifiers of the default imports.

Besides those of the figures, the class BasVal contains the following basic
values, which are bound to no identifier of the basis: the basic value
\kw{Elm.\allowbreak Kernel.\allowbreak Basics.\allowbreak log} of arity $1$
(Section~\ref{sec:appendix-dynamic-basis-elm}); the basic value
\kw{Task.command} of arity $1$, to which the environment of the texts of
\kw{Task} binds the hidden variable \kw{command}
(Section~\ref{sec:appendix-dynamic-basis-elm}), and whose $\APPLY$ is given in
Section~\ref{sec:effects-apply}; the basic values
\kw{Array.fromList} and \kw{Array.toList} of arity $1$, which
Chapter~\ref{ch:effects} applies to encode and decode arrays at ports
(Section~\ref{sec:appendix-dynamic-basis-arrays}); and the values of ports,
$\effoutport(\effp,\tau)$ and $\effinport(\effp,\tau)$, of arity $1$
(Section~\ref{sec:effects-ports}).

% ===========================================================================
\section{Values Written in Elm}\label{sec:appendix-dynamic-basis-elm}

A value of the basis whose definition in elm/core 1.0.5 is written in Elm is
defined by that text, except \kw{min}, \kw{max}, \kw{clamp} and
\kw{negate}, which are basic values (below); the other basic values are
those whose definition is a reference to the kernel. % A60, A48, A49
The texts of the values of \kw{Basics} marked ``Elm'' in
Figures~\ref{fig:dynamic-basis-basics} and~\ref{fig:dynamic-basis-operators}
are shown in Figure~\ref{fig:dynamic-basis-elm}, as they stand in the file
\kw{Basics.elm} of elm/core 1.0.5.

\begin{figure}[htbp]
\centering
\small
\begin{minipage}[t]{0.48\textwidth}
\footnotesize
\begin{verbatim}
abs n =
  if lt n 0 then -n else n

logBase base number =
  fdiv
    (Elm.Kernel.Basics.log number)
    (Elm.Kernel.Basics.log base)

radians angleInRadians =
  angleInRadians

degrees angleInDegrees =
  fdiv (mul angleInDegrees pi) 180

turns angleInTurns =
  mul (mul 2 pi) angleInTurns
\end{verbatim}
\end{minipage}\hfill
\begin{minipage}[t]{0.48\textwidth}
\footnotesize
\begin{verbatim}
fromPolar (radius, theta) =
  ( mul radius (cos theta)
  , mul radius (sin theta)
  )

toPolar ( x, y ) =
  ( sqrt (add (mul x x) (mul y y))
  , atan2 y x
  )

composeL g f x =
  g (f x)

composeR f g x =
  g (f x)

apR x f =
  f x

apL f x =
  f x

identity x =
  x

always a _ =
  a

never (JustOneMore nvr) =
  never nvr
\end{verbatim}
\end{minipage}
\caption{The Values of \kw{Basics} Written in Elm}\label{fig:dynamic-basis-elm}
\end{figure}

Each text of Figure~\ref{fig:dynamic-basis-elm} is a declaration $\dec$ that,
by the derived forms of Appendix~\ref{ch:appendix-derived-forms}, has the
form $\var\ \kw{=}\ \kw{\textbackslash}\,\atpat\ \kw{->}\ \expm$. Let
\[
  \begin{array}{@{}l@{}}
  E_{\mathrm{K}}=E_0+(\QE_{\mathrm{K}},\VE_{\mathrm{K}}),\quad\text{where}\\[0.5ex]
  \QE_{\mathrm{K}}=\{\kw{Elm.Kernel.Basics}\mapsto(\emptymap,\{\kw{log}\mapsto\kw{Elm.Kernel.Basics.log}\}),\\
  \phantom{\QE_{\mathrm{K}}=\{}\kw{Basics}\mapsto(\emptymap,\VE_{\mathrm{K}})\},
  \end{array}
\]
and $\VE_{\mathrm{K}}$ binds every identifier of \kw{Basics} that
Figures~\ref{fig:dynamic-basis-basics} and~\ref{fig:dynamic-basis-operators}
do not mark ``Elm'' to the value shown there, the hidden variables of the
operators included, and binds the hidden constructor \kw{JustOneMore} of the
closed type \kw{Never} to the constructor value \kw{JustOneMore} of arity
$1$. The qualifier \kw{Basics} is bound in $E_{\mathrm{K}}$ to the
environment of the identifiers of \kw{Basics} that $\VE_{\mathrm{K}}$ binds,
so that the long identifiers \kw{Basics.True}, \kw{Basics.False} and
\kw{Basics.negate} that the derived forms of
Appendix~\ref{ch:appendix-derived-forms} introduce into these texts (by
\kw{if} and by prefix minus) denote the values \kw{True}, \kw{False} and
\kw{negate} of Figure~\ref{fig:dynamic-basis-basics}; the texts are part of
the definition, not of a program, so the rule that a module cannot qualify
its own names (A38) does not apply to them.
% audit X18, V3; A38 governs programs only
No text refers
to a value that another text defines. The value bound to $\var$ in
$\ME_0(\kw{Basics})$ is $\VE(\var)$, where
\[
  \jdg{E_{\mathrm{K}}}{\dec}{\VE}
\]
by rule~\ruleref{r:d-dec-nonrec}, or, for \kw{never}, whose declaration is
recursive, by rule~\ruleref{r:d-dec-rec}. In each case $\VE$ is unique: it
binds $\var$ to a closure. % DECIDED: the kernel reference
% Elm.Kernel.Basics.log in the text of logBase denotes a basic value; kernel
% references are outside Standard Elm (A16), so the environment of the text
% supplies it. JustOneMore is bound only in that environment.
In the text of \kw{abs}, \kw{-n} is the derived form \kw{Basics.negate n};
the text of \kw{negate} itself is \kw{-n}, which by that derived form would
apply \kw{negate} to its own argument, so \kw{negate} is a basic value.
% DECIDED: negate is a basic value; its elm/core text is circular under the
% derived form of prefix minus.
\kw{min} and \kw{max} are basic values at every comparable type: at floats
their $\APPLY$ is the minimum and maximum of IEEE 754-2019, which their
elm/core texts do not give, and at every other comparable type it is the
result of those texts (Section~\ref{sec:appendix-dynamic-basis-order}). % A48
\kw{clamp} is a basic value likewise: at floats its $\APPLY$ is built on the
minimum and maximum of IEEE 754-2019, which its elm/core text does not
give, and at Int values it is the result of that text
(Sections~\ref{sec:appendix-dynamic-basis-floatarith}
and~\ref{sec:appendix-dynamic-basis-order}). % A48, A49
\kw{Debug.log}, \kw{Debug.todo} and \kw{Debug.toString} are kernel
references in elm/core, and are basic values.

The values of \kw{Platform.Cmd}, \kw{Platform.Sub} and \kw{Task} marked
``Elm'' in Figure~\ref{fig:dynamic-basis-others} are defined in the same way
by their texts in the files \kw{Platform/\allowbreak Cmd.elm},
\kw{Platform/\allowbreak Sub.elm} and \kw{Task.elm} of elm/core 1.0.5. The
text of a value of the module $\modid$ is evaluated in the environment
\[
  E_{\modid}=E_0+(\QE_{\modid},\VE_{\modid}),
\]
where $\QE_{\modid}$ binds each qualifier that the texts of $\modid$ use to
an environment binding the identifiers they write under it, to their values
in $\ME_0$ or, for a further function of elm/core, to the value of its text
(below), and $\VE_{\modid}$ binds the identifiers of the interfaces
of $\modid$ and of the modules it imports unqualified (among them the
constructors \kw{Ok} and \kw{Err}) to their values in $\ME_0$. The values
marked ``Elm'' are taken in an order in which each text refers only to values
already defined: \kw{Task.map} before \kw{Task.perform}, and \kw{Task.map2}
before \kw{Task.sequence}. For $\modid=\kw{Task}$, $\VE_{\kw{Task}}$ also binds
the two hidden identifiers of \kw{Task}
(Section~\ref{sec:appendix-static-basis-modenv}): the variable \kw{command}
to the basic value \kw{Task.command}, whose $\APPLY$ is that of
Figure~\ref{fig:effects-apply}, and the constructor \kw{Perform} of the
custom type \kw{MyCmd} to the constructor value \kw{Perform}, of arity $1$.
Neither is in $\ME_0(\kw{Task})$. The texts of \kw{Task.perform} and
\kw{Task.attempt} apply \kw{command} to a value \kw{Perform}~$\tk$, which
$\APPLY$ makes a leaf command of the task manager. % A60, A85; decision 5 (audit X9, Y3)
Where such a text refers to a further function of elm/core
(\kw{Task.sequence} applies \kw{List.foldr}), that function is defined by its
text likewise, and the kernel function \kw{Elm.Kernel.List.cons} is the list
constructor \kw{::}. Chapter~\ref{ch:effects} states what these texts give
(Section~\ref{sec:effects-apply}). % A60, A85

The texts are specialised over their \kw{number} variables with the rest of
the program (Section~\ref{sec:dynamic-core-numbers}): the copy of \kw{abs} used
at $\tFloat$ compares with the float $0$, and the constants \kw{180} and
\kw{2} in the texts of \kw{degrees} and \kw{turns}, whose type is $\tFloat$,
denote floats. % A87
Of the values of the basis written in Elm, only \kw{abs} has a type scheme
that quantifies a variable of class $\knum$
(Figure~\ref{fig:static-basis-basics-ve1}); the identifier \kw{abs} of
Figure~\ref{fig:dynamic-basis-basics} is bound to its default copy, at
$\tInt$. Besides this, the value environment of $\ME_0(\kw{Basics})$ binds
the new variable of the copy of the declaration of \kw{abs} at $\tFloat$
that the specialisation makes, to the value of that copy, evaluated in
$E_{\mathrm{K}}$ as the default copy is. An occurrence of \kw{abs} at
$\tFloat$ thus finds its copy: qualified, through the qualifier
\kw{Basics}; unqualified, through the default import
\kw{import Basics exposing (..)}, whose exposing list gives the whole of
that value environment; and through an exposing item \kw{abs} of an
import, which stands for the variable and the new variable of its copy
(Section~\ref{sec:dynamic-modules-reduced}).
So \kw{abs 1.5} is the float $1.5$. % A87; audit V1
% A49: clamp, whose type scheme also quantifies a variable of class number,
% is a basic value, so specialisation makes no copy of it; its APPLY
% dispatches on its arguments, as that of add does.

Thus the order in which a basis function written in Elm applies the functions
it is given, for instance \kw{List.foldr} from the last element of its list,
is fixed by its text. % A60
Consequences of the texts of Figure~\ref{fig:dynamic-basis-elm} that the
definitions of Section~\ref{sec:appendix-dynamic-basis-apply} make visible
are noted there. % A40, A60

% ===========================================================================
\section{The Function APPLY}\label{sec:appendix-dynamic-basis-apply}

For a basic value $b$ of arity $n$ and values $v_1,\ldots,v_n$,
$\APPLY(b,v_1\cdots v_n)$ is a finite set of values
(Section~\ref{sec:dynamic-core-basic}); in this section we separate the
arguments by commas, writing $\APPLY(b,v_1,\ldots,v_n)$. The set has exactly
one member except where stated: it has two members when an equality test
reaches two functions (Section~\ref{sec:appendix-dynamic-basis-eq}), and none
for \kw{Debug.todo} (Section~\ref{sec:appendix-dynamic-basis-debug}).
In this section $m$ and $n$ range over $\IntRange$, $x$, $y$ and $z$ over
$\effFlt$, and $s$ and $s'$ over $\effStr$.
The definitions cover every argument that the evaluation of a program that
elaborates supplies to $b$.

\textbf{Overloading.} The basic values \kw{add}, \kw{sub}, \kw{mul},
\kw{pow}, \kw{negate} and \kw{clamp} act on Int values and on floats; \kw{compare},
\kw{lt}, \kw{gt}, \kw{le}, \kw{ge}, \kw{min} and \kw{max} on all comparable
values; \kw{append} on
strings and on lists; \kw{eq} and \kw{neq} on all values. $\APPLY$ dispatches
on the special values it is given: each case is defined below for each kind of
argument. Before evaluation a program is specialised over its \kw{number}
variables, so that every integer constant has the type $\tInt$ or $\tFloat$,
and an integer constant of type $\tFloat$ denotes the nearest float
(Section~\ref{sec:dynamic-core-numbers}). % A87
Hence both arguments of a number operation are Int values or both are floats,
the three arguments of \kw{clamp} are all Int values or all floats, every
argument of a basic value of floats is a float, and an Int value never
meets a float in equality or order. % A49

% ===========================================================================
\subsection{Int Arithmetic}\label{sec:appendix-dynamic-basis-intarith}

For real $r$, $\lfloor r\rfloor$ is the greatest integer at most $r$,
$\lceil r\rceil$ the least integer at least $r$, and $\dbtrunc(r)$ is
$\lfloor r\rfloor$ when $r\geq 0$ and $\lceil r\rceil$ when $r<0$. On Int
values, $\APPLY$ is given in Figure~\ref{fig:dynamic-basis-int}. An exact
result is normalised into $\IntRange$; where an operation has no
mathematically correct result, the result is $0$. % D3, A40, A41, A42, A43

\begin{figure}[htbp]
\centering
\small
\[
\begin{array}{@{}r@{\ }c@{\ }l@{\quad}l@{}}
\APPLY(\kw{add},m,n) & = & \{\norm(m+n)\} & \\
\APPLY(\kw{sub},m,n) & = & \{\norm(m-n)\} & \\
\APPLY(\kw{mul},m,n) & = & \{\norm(m\cdot n)\} & \\
\APPLY(\kw{negate},n) & = & \{\norm(-n)\} & \\[1ex]
\APPLY(\kw{idiv},m,n) & = & \{\norm(\dbtrunc(m/n))\} & \text{if}\ n\neq 0\\
                      & = & \{0\} & \text{if}\ n=0\\[1ex]
\APPLY(\kw{modBy},n,m) & = & \{m-n\cdot\lfloor m/n\rfloor\} & \text{if}\ n\neq 0\\
                       & = & \{0\} & \text{if}\ n=0\\[1ex]
\APPLY(\kw{remainderBy},n,m) & = & \{m-n\cdot\dbtrunc(m/n)\} & \text{if}\ n\neq 0\\
                             & = & \{0\} & \text{if}\ n=0\\[1ex]
\APPLY(\kw{pow},m,n) & = & \{\norm(m^{n})\} & \text{if}\ n\geq 0\\
                     & = & \{1\} & \text{if}\ n<0\ \text{and}\ m=1\\
                     & = & \{(-1)^{n}\} & \text{if}\ n<0\ \text{and}\ m=-1\\
                     & = & \{0\} & \text{if}\ n<0\ \text{and}\ m\notin\{1,-1\}
\end{array}
\]
\caption{$\APPLY$ on Int Values}\label{fig:dynamic-basis-int}
\end{figure}

In the figure, $m^{0}=1$ for every $m$, including $0$. The first argument of
\kw{modBy} and \kw{remainderBy} is the divisor: \kw{modBy 4 -5} is $3$ and
\kw{remainderBy 4 -5} is $-1$. The following property, which may easily be
demonstrated, shows that \kw{modBy} and \kw{remainderBy} need no
normalisation.
\begin{property}
For $n\neq 0$, $m-n\cdot\lfloor m/n\rfloor$ is $0$ or has the sign of $n$,
$m-n\cdot\dbtrunc(m/n)$ is $0$ or has the sign of $m$, and both are less than
$|n|$ in absolute value; hence both are members of $\IntRange$.
\end{property}
Thus \kw{3 // 0}, \kw{modBy 0 3} and \kw{remainderBy 0 3} are all $0$, and
\kw{1 \textasciicircum{} -1} is $1$, \kw{(-1) \textasciicircum{} -3} is $-1$
and \kw{2 \textasciicircum{} -1} is $0$. % D3, A42, A43
The quotient \kw{($\min\IntRange$) // -1} is $\norm(-\min\IntRange)$. % A41
By the text of \kw{abs} (Figure~\ref{fig:dynamic-basis-elm}), \kw{abs n} is
$\norm(|n|)$. % A40

% ===========================================================================
\subsection{Floating-Point Arithmetic}\label{sec:appendix-dynamic-basis-floatarith}

On floats, \kw{add}, \kw{sub}, \kw{mul}, \kw{fdiv}, \kw{negate} and
\kw{sqrt} are the operations of IEEE 754-2019 on binary64 numbers, rounding
to nearest with ties to even; Figure~\ref{fig:dynamic-basis-float} gives them
case by case, the first applicable line giving the result.

\begin{figure}[htbp]
\centering
\small
\[
\begin{array}{@{}r@{\ }c@{\ }l@{\quad}l@{}}
\APPLY(\kw{negate},x) & = & \{\dbNaN\} & \text{if}\ x\ \text{is}\ \dbNaN\\
                      & = & \{(-1)\cdot x\} & \text{otherwise}\\[1ex]
\APPLY(\kw{add},x,y) & = & \{\dbNaN\} & \text{if}\ x\ \text{or}\ y\ \text{is}\ \dbNaN,\\
                     &   &             & \text{or}\ \{x,y\}=\{+\infty,-\infty\}\\
                     & = & \{x\} & \text{if}\ x\in\{+\infty,-\infty\}\\
                     & = & \{y\} & \text{if}\ y\in\{+\infty,-\infty\}\\
                     & = & \{-0\} & \text{if}\ x\ \text{and}\ y\ \text{are both}\ -0\\
                     & = & \{\dbfl(x+y)\} & \text{otherwise}\\[1ex]
\APPLY(\kw{sub},x,y) & = & \APPLY(\kw{add},x,z) & \text{where}\ \APPLY(\kw{negate},y)=\{z\}\\[1ex]
\APPLY(\kw{mul},x,y) & = & \{\dbNaN\} & \text{if}\ x\ \text{or}\ y\ \text{is}\ \dbNaN,\\
                     &   &             & \text{or one is a zero and one an infinity}\\
                     & = & \{\dbsgn(x)\dbsgn(y)\cdot\infty\} & \text{if}\ x\ \text{or}\ y\ \text{is an infinity}\\
                     & = & \{\dbsgn(x)\dbsgn(y)\cdot 0\} & \text{if}\ x\ \text{or}\ y\ \text{is a zero}\\
                     & = & \{\dbfl(x\cdot y)\} & \text{otherwise}\\[1ex]
\APPLY(\kw{fdiv},x,y) & = & \{\dbNaN\} & \text{if}\ x\ \text{or}\ y\ \text{is}\ \dbNaN,\\
                      &   &             & \text{or both are zeros, or both infinities}\\
                      & = & \{\dbsgn(x)\dbsgn(y)\cdot\infty\} & \text{if}\ x\ \text{is an infinity,}\ \text{or}\ y\ \text{a zero}\\
                      & = & \{\dbsgn(x)\dbsgn(y)\cdot 0\} & \text{if}\ x\ \text{is a zero,}\ \text{or}\ y\ \text{an infinity}\\
                      & = & \{\dbfl(x/y)\} & \text{otherwise}\\[1ex]
\APPLY(\kw{sqrt},x) & = & \{\dbNaN\} & \text{if}\ x\ \text{is}\ \dbNaN,\ \text{or}\ x<0\\
                    & = & \{x\} & \text{if}\ x\in\{+0,-0,+\infty\}\\
                    & = & \{\dbfl(\sqrt{x})\} & \text{otherwise}
\end{array}
\]
\[
\begin{array}{@{}r@{\ }c@{\ }l@{\quad}l@{}}
\APPLY(\kw{isNaN},x) & = & \{\kw{True}\} & \text{if}\ x\ \text{is}\ \dbNaN\\
                     & = & \{\kw{False}\} & \text{otherwise}\\[1ex]
\APPLY(\kw{isInfinite},x) & = & \{\kw{True}\} & \text{if}\ x\in\{+\infty,-\infty\}\\
                          & = & \{\kw{False}\} & \text{otherwise}
\end{array}
\]
\caption{$\APPLY$ on Floats: Arithmetic}\label{fig:dynamic-basis-float}
\end{figure}

Thus \kw{1 / 0} is $+\infty$, \kw{1 / -0} is $-\infty$, \kw{0 / 0} is NaN,
and \kw{x - x} is $+0$ for every finite \kw{x}. In the case $x<0$ of
\kw{sqrt}, $x$ is not a zero, so \kw{sqrt -0} is $-0$. By the text of
\kw{abs}, \kw{abs -0} is $-0$ and \kw{abs} of NaN is NaN, since
\kw{lt} is false for both (Section~\ref{sec:appendix-dynamic-basis-order}).

At floats, \kw{min} and \kw{max} are the operations \emph{minimum} and
\emph{maximum} of IEEE 754-2019, the first applicable line giving the
result: % A48
\[
\begin{array}{@{}r@{\ }c@{\ }l@{\quad}l@{}}
\APPLY(\kw{min},x,y) & = & \{\dbNaN\} & \text{if}\ x\ \text{or}\ y\ \text{is}\ \dbNaN\\
                     & = & \{-0\} & \text{if}\ x\ \text{and}\ y\ \text{are zeros and one of them is}\ {-0}\\
                     & = & \{x\} & \text{if}\ x\leq y\\
                     & = & \{y\} & \text{otherwise}\\[1ex]
\APPLY(\kw{max},x,y) & = & \{\dbNaN\} & \text{if}\ x\ \text{or}\ y\ \text{is}\ \dbNaN\\
                     & = & \{+0\} & \text{if}\ x\ \text{and}\ y\ \text{are zeros and one of them is}\ {+0}\\
                     & = & \{x\} & \text{if}\ x\geq y\\
                     & = & \{y\} & \text{otherwise}.
\end{array}
\]
These lines apply when both arguments are floats; at the other comparable
values $\APPLY$ takes the cases of
Section~\ref{sec:appendix-dynamic-basis-order}. The case is selected by the
arguments, as for \kw{lt}: the type schemes of \kw{min} and \kw{max}
quantify a variable of class $\kcomp$, not $\knum$, so number
specialisation makes no copy of them
(Section~\ref{sec:dynamic-core-numbers}). The result does not depend
on the order of the arguments: with \kw{nan} for \kw{0/0}, \kw{max nan 1}
and \kw{max 1 nan} are both NaN, \kw{max 0 -0} and \kw{max -0 0} are $+0$,
and \kw{min -0 0} and \kw{min 0 -0} are $-0$. % A48

At floats, \kw{clamp} is built on these operations: % A48, A49
\[
\begin{array}{@{}r@{\ }c@{\ }l@{\quad}l@{}}
\APPLY(\kw{clamp},x,y,z) & = & \APPLY(\kw{min},z',y) &\\
 & & \multicolumn{2}{@{}l@{}}{\quad\text{where}\ \APPLY(\kw{max},z,x)=\{z'\},}
\end{array}
\]
when $x$, $y$ and $z$ are floats; the first argument is the lower bound,
the second the upper bound, and \kw{clamp lo hi x} is
\kw{min (max x lo) hi}. At Int values $\APPLY$ takes the cases of
Section~\ref{sec:appendix-dynamic-basis-order}; the case is selected by the
arguments, as for \kw{add}. When \kw{lo} is greater than \kw{hi}, the
result is \kw{hi} unless an argument is NaN: \kw{clamp 10 0 5} is $+0$.
A NaN in any argument gives NaN: \kw{clamp 0 1 nan}, \kw{clamp nan 1 0.5}
and \kw{clamp 0 nan 0.5} are all NaN. Zeros are ordered as by
\kw{max} and \kw{min}: \kw{clamp 0 1 -0} is $+0$, the maximum of $-0$ and
$+0$ being $+0$, and \kw{clamp -1 -0 0} is $-0$.
% A48, A49

% ===========================================================================
\subsection{Powers and Elementary Functions}\label{sec:appendix-dynamic-basis-elementary}

The power at floats, the trigonometric functions, and the logarithm
\kw{Elm.\allowbreak Kernel.\allowbreak Basics.\allowbreak log} on which \kw{logBase} is built, denote the real
functions $x^{y}$, $\sin$, $\cos$, $\tan$, $\arcsin$, $\arccos$, $\arctan$,
the natural logarithm $\ln$, and the two-argument arctangent, rounded to a
float, with the special values that IEEE 754-2019 (Section 9.2) recommends.
% A65
The two-argument arctangent of a point $(x,y)\neq(0,0)$ is the angle
$\theta$ with $-\dbpinum<\theta\leq\dbpinum$ such that
$(x,y)=\rho\cdot(\cos\theta,\sin\theta)$ for some $\rho>0$; the first argument
of \kw{atan2} is $y$. The power is given in
Figure~\ref{fig:dynamic-basis-pow} and the other functions in
Figure~\ref{fig:dynamic-basis-elem}; in each, the first applicable line gives
the result. An integer $k$ is \emph{odd} when $k=2j+1$ for an integer $j$.

\begin{figure}[htbp]
\centering
\small
\[
\begin{array}{@{}r@{\ }c@{\ }l@{\quad}l@{}}
\APPLY(\kw{pow},x,y) & = & \{1\} & \text{if}\ y\ \text{is a zero}\\
  & = & \{1\} & \text{if}\ x\ \text{is}\ 1\\
  & = & \{\dbNaN\} & \text{if}\ x\ \text{or}\ y\ \text{is}\ \dbNaN\\
  & = & \{\dbsgn(x)\cdot\infty\} & \text{if}\ x\ \text{is a zero and}\ y<0\ \text{an odd integer}\\
  & = & \{+\infty\} & \text{if}\ x\ \text{is a zero and}\ y<0\\
  & = & \{x\} & \text{if}\ x\ \text{is a zero and}\ y>0\ \text{an odd integer}\\
  & = & \{+0\} & \text{if}\ x\ \text{is a zero and}\ y>0\\
  & = & \{1\} & \text{if}\ x=-1\ \text{and}\ y\in\{+\infty,-\infty\}\\
  & = & \{+0\} & \text{if}\ |x|<1\ \text{and}\ y=+\infty\\
  & = & \{+0\} & \text{if}\ |x|>1\ \text{and}\ y=-\infty\\
  & = & \{+\infty\} & \text{if}\ |x|>1\ \text{and}\ y=+\infty\\
  & = & \{+\infty\} & \text{if}\ |x|<1\ \text{and}\ y=-\infty\\
  & = & \{+0\} & \text{if}\ x=+\infty\ \text{and}\ y<0\\
  & = & \{+\infty\} & \text{if}\ x=+\infty\ \text{and}\ y>0\\
  & = & \{-0\} & \text{if}\ x=-\infty\ \text{and}\ y<0\ \text{an odd integer}\\
  & = & \{-\infty\} & \text{if}\ x=-\infty\ \text{and}\ y>0\ \text{an odd integer}\\
  & = & \{+0\} & \text{if}\ x=-\infty\ \text{and}\ y<0\\
  & = & \{+\infty\} & \text{if}\ x=-\infty\ \text{and}\ y>0\\
  & = & \{\dbNaN\} & \text{if}\ x<0\ \text{and}\ y\ \text{is not an integer}\\
  & = & \{\dbfl(x^{y})\} & \text{otherwise}
\end{array}
\]
\caption{$\APPLY$ on Floats: Power}\label{fig:dynamic-basis-pow}
\end{figure}

The first two lines apply even when the other argument is NaN. In the last
line of Figure~\ref{fig:dynamic-basis-pow}, $x$ and $y$ are finite
and $x$ is not a zero; when $x<0$, $y$ is an integer and
$x^{y}=(-1)^{y}|x|^{y}$. Thus \kw{1 \textasciicircum{} (1/0)} is $1$, and
\kw{(-8) \textasciicircum{} (1/3)} is NaN. % A65

\begin{figure}[htbp]
\centering
\small
\[
\begin{array}{@{}r@{\ }c@{\ }l@{\quad}l@{}}
\APPLY(b,\dbNaN) & = & \{\dbNaN\} & \text{for each}\ b\ \text{of the figure}\\[1ex]
\APPLY(b,x) & = & \{x\} & \text{if}\ x\ \text{is a zero, for}\ b\in\{\kw{sin},\kw{tan},\kw{asin},\kw{atan}\}\\
\APPLY(b,x) & = & \{\dbNaN\} & \text{if}\ x\in\{+\infty,-\infty\},\ \text{for}\ b\in\{\kw{sin},\kw{cos},\kw{tan}\}\\
\APPLY(b,x) & = & \{\dbNaN\} & \text{if}\ |x|>1,\ \text{for}\ b\in\{\kw{asin},\kw{acos}\}\\
\APPLY(\kw{atan},x) & = & \{\dbsgn(x)\cdot\dbfl(\dbpinum/2)\} & \text{if}\ x\in\{+\infty,-\infty\}
\end{array}
\]
\[
\begin{array}{@{}r@{\ }c@{\ }l@{\quad}l@{}}
\APPLY(\kw{Elm.Kernel.Basics.log},x) & = & \{-\infty\} & \text{if}\ x\ \text{is a zero}\\
  & = & \{\dbNaN\} & \text{if}\ x<0\\
  & = & \{+\infty\} & \text{if}\ x=+\infty\\[1ex]
\APPLY(b,x) & = & \{\dbfl(f(x))\} & \text{otherwise, for each}\ b
\end{array}
\]
\[
\begin{array}{@{}r@{\ }c@{\ }l@{\quad}l@{}}
\APPLY(\kw{atan2},y,x) & = & \{\dbNaN\} & \text{if}\ x\ \text{or}\ y\ \text{is}\ \dbNaN\\
  & = & \{y\} & \text{if}\ y\ \text{is a zero, and}\ x\ \text{is}\ +0\ \text{or}\ x>0\\
  & = & \{\dbsgn(y)\cdot\dbfl(\dbpinum)\} & \text{if}\ y\ \text{is a zero, and}\ x\ \text{is}\ -0\ \text{or}\ x<0\\
  & = & \{\dbsgn(y)\cdot\dbfl(\dbpinum/2)\} & \text{if}\ x\ \text{is a zero}\\
  & = & \{\dbsgn(y)\cdot\dbfl(\dbpinum/2)\} & \text{if}\ y\ \text{is an infinity and}\ x\ \text{is finite}\\
  & = & \{\dbsgn(y)\cdot\dbfl(3\dbpinum/4)\} & \text{if}\ y\ \text{is an infinity and}\ x=-\infty\\
  & = & \{\dbsgn(y)\cdot\dbfl(\dbpinum/4)\} & \text{if}\ y\ \text{is an infinity and}\ x=+\infty\\
  & = & \{\dbsgn(y)\cdot\dbfl(\dbpinum)\} & \text{if}\ x=-\infty\\
  & = & \{\dbsgn(y)\cdot 0\} & \text{if}\ x=+\infty\\
  & = & \{\dbfl(\theta)\} & \text{otherwise}
\end{array}
\]
\caption{$\APPLY$ on Floats: Elementary Functions}\label{fig:dynamic-basis-elem}
\end{figure}

In the first part of Figure~\ref{fig:dynamic-basis-elem}, $b$ ranges over
\kw{sin}, \kw{cos}, \kw{tan}, \kw{asin}, \kw{acos}, \kw{atan} and
\kw{Elm.\allowbreak Kernel.\allowbreak Basics.\allowbreak log}, and $f$ is the
real function of $b$: $\sin$, $\cos$, $\tan$, $\arcsin$, $\arccos$, $\arctan$
or $\ln$; on the arguments of the line $\{\dbfl(f(x))\}$, $f(x)$ is a real
number. In the last line for \kw{atan2}, $\theta$ is the two-argument
arctangent of $(x,y)$. The values \kw{e} and \kw{pi} are $\dbfl(\dbeuler)$
and $\dbfl(\dbpinum)$, and \kw{logBase}, \kw{degrees}, \kw{radians},
\kw{turns}, \kw{fromPolar} and \kw{toPolar} are given by their texts
(Figure~\ref{fig:dynamic-basis-elm}), so that \kw{logBase 10 1000} is the
quotient of two rounded logarithms, rounded again.

The lines $\{\dbfl(x^{y})\}$, $\{\dbfl(f(x))\}$ and $\{\dbfl(\theta)\}$ of
Figures~\ref{fig:dynamic-basis-pow} and~\ref{fig:dynamic-basis-elem} give the
correctly rounded result. On those arguments, and only there, the accuracy
with which an implementation computes the result is implementation-dependent,
and an implementation may give a float near the real result in its place
(Section~\ref{sec:appendix-implementation-float}). The operations of
Figure~\ref{fig:dynamic-basis-float} allow no such latitude. % A65, D8

% ===========================================================================
\subsection{Conversions}\label{sec:appendix-dynamic-basis-conv}

An Int value becomes a float by rounding, and a float becomes an Int value by
\emph{saturation} into $\IntRange$:
\[
  \dbsat(k)=\left\{
  \begin{array}{@{}ll@{}}
  k & \text{if}\ k\in\IntRange\\
  \max\IntRange & \text{if}\ k=+\infty\ \text{or}\ k>\max\IntRange\\
  \min\IntRange & \text{if}\ k=-\infty\ \text{or}\ k<\min\IntRange
  \end{array}\right.
\]
for $k\in\mathbb{Z}\cup\{+\infty,-\infty\}$. The conversions are given in Figure~\ref{fig:dynamic-basis-conv}. % A44, A45

\begin{figure}[htbp]
\centering
\small
\[
\begin{array}{@{}r@{\ }c@{\ }l@{\quad}l@{}}
\APPLY(\kw{toFloat},n) & = & \{\dbfl(n)\} & \\[1ex]
\APPLY(\kw{Char.toCode},\dbch) & = & \{\dbcode(\dbch)\} & \\
\APPLY(\kw{Char.fromCode},n) & = & \{\dbch\} & \text{if}\ n\ \text{is a Unicode scalar value}\\
                              &   &            & \text{and}\ \dbcode(\dbch)=n\\
                              & = & \{\mathit{other}(n)\} & \text{otherwise}\\[1ex]
\APPLY(\kw{round},x) & = & \{\dbsat(\lfloor x+\tfrac{1}{2}\rfloor)\} & \text{if}\ x\in\effFltFin\\
\APPLY(\kw{floor},x) & = & \{\dbsat(\lfloor x\rfloor)\} & \text{if}\ x\in\effFltFin\\
\APPLY(\kw{ceiling},x) & = & \{\dbsat(\lceil x\rceil)\} & \text{if}\ x\in\effFltFin\\
\APPLY(\kw{truncate},x) & = & \{\dbsat(\dbtrunc(x))\} & \text{if}\ x\in\effFltFin\\[1ex]
\APPLY(b,x) & = & \{\dbsat(x)\} & \text{if}\ x\in\{+\infty,-\infty\}\\
\APPLY(b,\dbNaN) & = & \{0\} &
\end{array}
\]
\caption{$\APPLY$ on Conversions}\label{fig:dynamic-basis-conv}
\end{figure}

In the last two lines $b$ is any of \kw{round}, \kw{floor}, \kw{ceiling} and
\kw{truncate}. The sum $x+\frac{1}{2}$ is exact, not rounded to a float, so
\kw{round} rounds a half upward: \kw{round 2.5} is $3$ and \kw{round -2.5} is
$-2$. A float zero of either sign converts to the one Int zero, and
\kw{toFloat 0} is $+0$, so \kw{1 / toFloat (negate 0)} is $+\infty$. % A45

% ===========================================================================
\subsection{Equality}\label{sec:appendix-dynamic-basis-eq}

The \emph{function values} are the closures, the basic values with their
partial applications, and the constructor values awaiting arguments:
\[
  \dbFun=\Closure\cup\BasApp\cup\{\con\ v_1\cdots v_k\in\ConVal\mid k<\arity(\con)\}.
\]
Equality is structural on all values, and is defined by a set
$\dbEq(v,v')\subseteq\{\kw{True},\kw{False}\}$ in
Figure~\ref{fig:dynamic-basis-eq}, by recursion on the structure of $v$ and
$v'$ (values are finite, and the recursion does not enter a function value).
Then
\[
  \begin{array}{rcl}
  \APPLY(\kw{eq},v,v') &=& \dbEq(v,v'),\\
  \APPLY(\kw{neq},v,v') &=& \{\kw{True}\mid\kw{False}\in\dbEq(v,v')\}\\
  &&{}\cup\{\kw{False}\mid\kw{True}\in\dbEq(v,v')\}.
  \end{array}
\]
% A35, A50, A51

\begin{figure}[htbp]
\centering
\small
\begin{enumerate}
\item If $v,v'\in\dbFun$, then $\dbEq(v,v')=\{\kw{True},\kw{False}\}$.
\item If $v$ and $v'$ are floats: $\dbEq(v,v')=\{\kw{True}\}$ if neither is
  NaN and $v=v'$ (so that $+0$ equals $-0$), and $\{\kw{False}\}$ otherwise.
\item If $v$ and $v'$ are both Int values, both characters, both strings,
  both the unit value, both process identifiers, or both boundary values
  (Chapter~\ref{ch:effects}) that are $\effnull$, $\efftrue$, $\efffalse$,
  numbers, strings or byte sequences:
  $\dbEq(v,v')=\{\kw{True}\}$ if $v$ and $v'$ are the same, and
  $\{\kw{False}\}$ otherwise.
\item If $v$ and $v'$ have the same form, with components $v_1,\ldots,v_k$ and
  $v'_1,\ldots,v'_k$ ($k\geq 0$), as listed below, then
  $\kw{True}\in\dbEq(v,v')$ if $\kw{True}\in\dbEq(v_i,v'_i)$ for every $i$, and
  $\kw{False}\in\dbEq(v,v')$ if $\kw{False}\in\dbEq(v_i,v'_i)$ for some $i$.
  The forms are:
  \begin{itemize}
  \item tuples with the same number of components;
  \item constructor values $\con\ v_1\cdots v_k$ and $\con\ v'_1\cdots v'_k$
    with the same constructor name and $k=\arity(\con)$;
  \item records $r$ and $r'$ with $\Dom r=\Dom r'$, the components being
    $r(\lab)$ and $r'(\lab)$ for each $\lab\in\Dom r$;
  \item arrays of boundary values (JSON arrays among them) with the same
    number of elements, and objects of boundary values with the same set of
    keys, the components being the values at each key; % A92
  \item effect values (Chapter~\ref{ch:effects}) of the same form, with the
    same primitives, managers and port names, the components being the values,
    boundary values, tasks, commands and subscriptions they contain.
  \end{itemize}
\item Otherwise $\dbEq(v,v')=\{\kw{False}\}$.
\end{enumerate}
\caption{Equality of Values}\label{fig:dynamic-basis-eq}
\end{figure}

The clauses are tried in order. Lists, \kw{Maybe} and \kw{Result} values,
\kw{Bool} and \kw{Order} values, arrays and the values of type \kw{Bytes}
are constructor values, so they are compared by the fourth clause: two lists
are equal when they have the same length and equal elements. A record type gives all its records the same
domain, so the condition $\Dom r=\Dom r'$ always holds in a program that
elaborates. The following property may easily be demonstrated.
\begin{property}
$\dbEq(v,v')$ is non-empty. It is $\{\kw{True},\kw{False}\}$ only if, at some
depth, corresponding components of $v$ and $v'$ are both function values, and
it is $\{\kw{False}\}$ whenever some pair of corresponding components that are
not function values is unequal.
\end{property}
Equality is thus deterministic except when the comparison reaches two
function values, where it may give either \kw{True} or \kw{False}:
\kw{(\textbackslash x -> x) == (\textbackslash x -> x)} may evaluate to
either. This is the only non-determinism in the dynamic semantics of the
Core (rule~\ruleref{r:d-app-basic-apply}). % A51, R9
Because equality looks only at structure, the identity of values plays no
part. NaN is unequal to every float, itself included, so a value with a NaN
component is not equal to itself: if \kw{t} is \kw{(0/0, 1)}, then
\kw{t == t} is \kw{False}. % A50
The two zeros of the floats are equal. Implementations whose equality
consults identity are recorded in
Section~\ref{sec:appendix-implementation-equality}.

% ===========================================================================
\subsection{Order}\label{sec:appendix-dynamic-basis-order}

The \emph{comparable values} are the Int values, the floats, the characters,
the strings, and the lists and tuples whose elements or components are
comparable values; they are the values of the types admissible for the class
$\kcomp$ (Section~\ref{sec:static-core-classes}). For comparable values $v$ and
$v'$ of one type, $\dbcmp(v,v')\in\{\kw{LT},\kw{EQ},\kw{GT}\}$ is defined in
Figure~\ref{fig:dynamic-basis-order}. Then
\[
  \APPLY(\kw{compare},v,v')=\{\dbcmp(v,v')\}.
\]
% A46, A52, A53

\begin{figure}[htbp]
\centering
\small
\begin{enumerate}
\item Int values $m$, $n$: $\dbcmp(m,n)$ is \kw{LT} if $m<n$, \kw{EQ} if
  $m=n$, and \kw{GT} if $m>n$.
\item Floats $x$, $y$: $\dbcmp(x,y)$ is \kw{EQ} if both are NaN, \kw{GT} if only $x$ is
  NaN, \kw{LT} if only $y$ is NaN, and otherwise \kw{LT}, \kw{EQ} or \kw{GT}
  as $x<y$, $x=y$ or $x>y$.
\item Characters: $\dbcmp(\dbch,\dbch')=\dbcmp(\dbcode(\dbch),\dbcode(\dbch'))$.
\item Strings $s=\dbch_1\cdots\dbch_k$ and $s'=\dbch'_1\cdots\dbch'_l$: if
  $\dbch_i\neq\dbch'_i$ for some $i\leq\min(k,l)$, and $i$ is the least such,
  then $\dbcmp(s,s')=\dbcmp(\dbch_i,\dbch'_i)$; otherwise
  $\dbcmp(s,s')=\dbcmp(k,l)$.
\item Lists: $\dbcmp(\kw{[]},\kw{[]})=\kw{EQ}$,
  $\dbcmp(\kw{[]},v_1\mathbin{\kw{::}}v_2)=\kw{LT}$,
  $\dbcmp(v_1\mathbin{\kw{::}}v_2,\kw{[]})=\kw{GT}$, and
  $\dbcmp(v_1\mathbin{\kw{::}}v_2,\,v'_1\mathbin{\kw{::}}v'_2)$ is
  $\dbcmp(v_1,v'_1)$ if that is not \kw{EQ}, and $\dbcmp(v_2,v'_2)$ otherwise.
\item Tuples $(v_1,\ldots,v_k)$ and $(v'_1,\ldots,v'_k)$: $\dbcmp(v_i,v'_i)$
  for the least $i$ with $\dbcmp(v_i,v'_i)\neq\kw{EQ}$, and \kw{EQ} if there is
  none.
\end{enumerate}
\caption{Order of Comparable Values}\label{fig:dynamic-basis-order}
\end{figure}

\begin{property}
For comparable values $v$, $v'$ and $v''$ of one type,
$\dbcmp(v,v')=\kw{LT}$ if and only if $\dbcmp(v',v)=\kw{GT}$,
$\dbcmp(v,v)=\kw{EQ}$, and if $\dbcmp(v,v')\neq\kw{GT}$ and
$\dbcmp(v',v'')\neq\kw{GT}$ then $\dbcmp(v,v'')\neq\kw{GT}$.
\end{property}
Thus \kw{compare} is a total order on each comparable type, up to the values
it makes \kw{EQ}, and sorting by it has a defined result. It is
\kw{EQ} for the two zeros and for two NaNs, and puts NaN above every other
float; it therefore disagrees with equality at NaN, which is unequal to
itself. % A46
Characters are ordered by $\dbcode$, and strings by their sequences of
characters, so the order of one-character strings is that of their
characters. % A53

The order operators are the IEEE 754 predicates at floats and are derived
from $\dbcmp$ at every other comparable type: % A47
\[
  \begin{array}{@{}r@{\ }c@{\ }l@{\quad}l@{}}
  \APPLY(\kw{lt},x,y) & = & \{\kw{True}\} & \text{if neither is NaN and}\ x<y\\
                      & = & \{\kw{False}\} & \text{otherwise}\\
  \APPLY(\kw{le},x,y) & = & \{\kw{True}\} & \text{if neither is NaN and}\ x\leq y\\
                      & = & \{\kw{False}\} & \text{otherwise}\\
  \APPLY(\kw{gt},x,y) & = & \APPLY(\kw{lt},y,x) & \\
  \APPLY(\kw{ge},x,y) & = & \APPLY(\kw{le},y,x) & \\[1ex]
  \APPLY(\kw{lt},v,v') & = & \{\kw{True}\} & \text{if}\ \dbcmp(v,v')=\kw{LT}\\
                       & = & \{\kw{False}\} & \text{otherwise}\\
  \APPLY(\kw{le},v,v') & = & \{\kw{True}\} & \text{if}\ \dbcmp(v,v')\neq\kw{GT}\\
                       & = & \{\kw{False}\} & \text{otherwise}\\
  \APPLY(\kw{gt},v,v') & = & \APPLY(\kw{lt},v',v) & \\
  \APPLY(\kw{ge},v,v') & = & \APPLY(\kw{le},v',v) &
  \end{array}
\]
where in the first six lines $x$ and $y$ are floats, and in the last six $v$ and
$v'$ are comparable values of any other kind. Hence \kw{nan < 1},
\kw{nan > 1} and \kw{nan >= nan} are \kw{False}, where \kw{nan} is
\kw{0/0}, while \kw{[ nan ] > [ 1 ]} is \kw{True}, the list being compared by
$\dbcmp$. % A47

At comparable values other than floats, \kw{min} and \kw{max} give the
results of their elm/core texts, \kw{min x y = if lt x y then x else y} and
\kw{max x y = if gt x y then x else y}: % A48
\[
  \begin{array}{@{}r@{\ }c@{\ }l@{\quad}l@{}}
  \APPLY(\kw{min},v,v') & = & \{v\} & \text{if}\ \dbcmp(v,v')=\kw{LT}\\
                        & = & \{v'\} & \text{otherwise}\\
  \APPLY(\kw{max},v,v') & = & \{v\} & \text{if}\ \dbcmp(v,v')=\kw{GT}\\
                        & = & \{v'\} & \text{otherwise},
  \end{array}
\]
where $v$ and $v'$ are comparable values that are not floats; the cases at
floats are those of Section~\ref{sec:appendix-dynamic-basis-floatarith}.
Thus \kw{max [ nan ] [ 1 ]} and \kw{max [ 1 ] [ nan ]} are both
\kw{[ nan ]}, the lists being compared by $\dbcmp$, which puts NaN above
every other float. % A48

At Int values, \kw{clamp} gives the result of its elm/core text,
\[
\begin{array}{@{}l@{}}
\kw{clamp low high number =}\\
\quad\kw{if lt number low then low}\\
\quad\kw{else if gt number high then high else number},
\end{array}
\]
the first applicable line giving the result: % A49
\[
  \begin{array}{@{}r@{\ }c@{\ }l@{\quad}l@{}}
  \APPLY(\kw{clamp},v,v',v'') & = & \{v\} & \text{if}\ \APPLY(\kw{lt},v'',v)=\{\kw{True}\}\\
                              & = & \{v'\} & \text{if}\ \APPLY(\kw{gt},v'',v')=\{\kw{True}\}\\
                              & = & \{v''\} & \text{otherwise},
  \end{array}
\]
where $v$, $v'$ and $v''$ are not floats, and are therefore Int values, the
type scheme of \kw{clamp} quantifying a variable of class $\knum$; the case
at floats is that of Section~\ref{sec:appendix-dynamic-basis-floatarith}.
When \kw{lo} is greater than \kw{hi}, \kw{clamp lo hi x} at Int values is
\kw{lo} if \kw{x < lo}, else \kw{hi} if \kw{x > hi}, else \kw{x}:
\kw{clamp 10 0 5} is $10$, where at floats it is $+0$. % A49

% ===========================================================================
\subsection{Booleans, Appending and Functions}\label{sec:appendix-dynamic-basis-bool}

The Boolean values are the constructor values \kw{True} and \kw{False}.
\[
  \begin{array}{@{}r@{\ }c@{\ }l@{\quad}l@{}}
  \APPLY(\kw{not},v) & = & \{\kw{False}\} & \text{if}\ v=\kw{True}\\
                     & = & \{\kw{True}\} & \text{otherwise}\\
  \APPLY(\kw{and},v,v') & = & \{\kw{True}\} & \text{if}\ v=v'=\kw{True}\\
                        & = & \{\kw{False}\} & \text{otherwise}\\
  \APPLY(\kw{or},v,v') & = & \{\kw{False}\} & \text{if}\ v=v'=\kw{False}\\
                       & = & \{\kw{True}\} & \text{otherwise}\\
  \APPLY(\kw{xor},v,v') & = & \{\kw{True}\} & \text{if}\ v\neq v'\\
                        & = & \{\kw{False}\} & \text{otherwise}
  \end{array}
\]
The basic values \kw{and} and \kw{or} receive both arguments as values, so the
operator values \kw{(\&\&)} and \kw{(||)} are strict; only the infixed forms
\kw{e1 \&\& e2} and \kw{e1 || e2}, which are derived forms, evaluate their
second operand conditionally. % A56

The basic value \kw{append} concatenates strings and lists; for strings,
$s\,s'$ is the sequence of the characters of $s$ followed by those of $s'$:
\[
  \begin{array}{@{}r@{\ }c@{\ }l@{}}
  \APPLY(\kw{append},s,s') & = & \{s\,s'\}\\
  \APPLY(\kw{append},\kw{[]},v) & = & \{v\}\\
  \APPLY(\kw{append},v_1\mathbin{\kw{::}}v_2,v) & = & \{v_1\mathbin{\kw{::}}v'\}
  \end{array}
\]
where, in the last line, $\APPLY(\kw{append},v_2,v)=\{v'\}$.
The composition and application operators and \kw{identity}, \kw{always} and
\kw{never} are given by their texts (Figure~\ref{fig:dynamic-basis-elm}). No
value has the type \kw{Never}, so \kw{never} is never applied.

% ===========================================================================
\subsection{Debugging}\label{sec:appendix-dynamic-basis-debug}

\[
  \begin{array}{@{}r@{\ }c@{\ }l@{}}
  \APPLY(\kw{Debug.log},s,v) & = & \{v\}\\
  \APPLY(\kw{Debug.todo},s) & = & \emptyset\\
  \APPLY(\kw{Debug.toString},v) & = & \{s_v\}
  \end{array}
\]
where $s_v$ is the string that the implementation gives for $v$.
Thus \kw{Debug.log} returns its second argument; any output it makes lies outside
the definition. \kw{Debug.todo} has no result, so an application of it has no
value (Section~\ref{sec:dynamic-core-basic}). The string of
\kw{Debug.toString} is determined by its argument, in a way that is
implementation-dependent (Section~\ref{sec:appendix-implementation-order}).
% A62, D8

% ===========================================================================
\subsection{Arrays}\label{sec:appendix-dynamic-basis-arrays}

We take the class Con to contain, besides \kw{[]} and \kw{::}, a third
constructor name $\dbarr$ that no identifier spells, with
$\arity(\dbarr)=1$. The values of the type \kw{Array} are the constructor
values $\dbarr\ v$ with $v$ a list: the array of the elements of $v$, in
order. The two basic values that Chapter~\ref{ch:effects} applies to arrays
are
\[
  \begin{array}{@{}r@{\ }c@{\ }l@{}}
  \APPLY(\kw{Array.fromList},v) & = & \{\dbarr\ v\},\\
  \APPLY(\kw{Array.toList},\dbarr\ v) & = & \{v\}.
  \end{array}
\]
Two arrays are equal when their lists are, by the fourth clause of
Figure~\ref{fig:dynamic-basis-eq}. The other functions of the library module
\kw{Array} lie outside the definition (Appendix~\ref{ch:appendix-static-basis}).
% DECIDED: elm/core writes Array.fromList and Array.toList in Elm over the
% kernel type JsArray; kernel code being outside the definition, the values
% of Array are given abstractly here, and these two functions agree with that
% text: toList (fromList l) is l.

\subsection{Bytes}\label{sec:appendix-dynamic-basis-bytes}

We take the class Con to contain also a fourth constructor name $\dbbin$ that
no identifier spells, with $\arity(\dbbin)=1$. The values of the type
\kw{Bytes} are the constructor values $\dbbin\ v$ with $v$ a list of Int
values each from $0$ to $255$; the value $\dbbin\,[o_1,\ldots,o_k]$ is the
finite sequence of the bytes $o_1,\ldots,o_k$, in order. % A92
Two values of type \kw{Bytes} are equal exactly when their sequences are, by
the fourth clause of Figure~\ref{fig:dynamic-basis-eq}: equality on them is
structural, like that on lists. % A50
The values of \kw{Bytes} meet the definition only at ports and in flags,
where Chapter~\ref{ch:effects} relates them to the byte sequences of the
boundary (Section~\ref{sec:effects-ports}); the functions of the library
\kw{elm/bytes} lie outside the definition
(Appendix~\ref{ch:appendix-static-basis}).
% DECIDED (A92): elm/bytes represents Bytes by a kernel DataView; kernel code
% being outside the definition, the values are given abstractly here, as for
% Array, and the interface of Bytes in ME0 is empty.

% ===========================================================================
\subsection{Effect Values}\label{sec:appendix-dynamic-basis-effects}

For the basic values of \kw{Task}, \kw{Process}, \kw{Platform.Cmd},
\kw{Platform.Sub} and \kw{Platform}, and for \kw{Task.command}, $\APPLY$ and
the arities are given in Figure~\ref{fig:effects-apply}
(Section~\ref{sec:effects-apply}); for the values of ports,
$\effoutport(\effp,\tau)$ and $\effinport(\effp,\tau)$, in
Section~\ref{sec:effects-ports}. % D10, D11
